%!TEX root = ../LLMSeminar.tex
% ──────────────────────────────────────────────────────────────
%  Section 5 — Applications and Q\&A Highlights
% ──────────────────────────────────────────────────────────────

\section{Applications and discussion}\label{sec:applications}

\subsection{What you can do with infinite cognition}

If you have effectively infinite cognition at your disposal, what does
that enable?  Here are several examples demonstrated during the talk,
all produced autonomously by Claude Code:

\begin{enumerate}[(1)]
  \item \textbf{Meeting transcription and enrichment.}  A Zoom recording
    of a blackboard discussion was processed end-to-end: Claude
    transcribed the audio via Whisper, extracted screenshots of the
    board from the video, recognised barely readable equations, wrote
    \LaTeX\ of everything it saw, and enriched the notes with the
    discussion context.  Fully automated, not a word typed by hand.

  \item \textbf{Custom presentation software.}  A colleague wanted
    presentation software that used Typst.  It did not exist, so he
    asked Claude to write it.  In a weekend he had working software.

  \item \textbf{Auto-formalisation.}  Mathematical results from research
    papers formalised into Lean~4 proofs, checked by the compiler.

  \item \textbf{Physically correct ray tracing.}  A ray tracer that
    correctly represents the spontaneous emission of an excited hydrogen
    atom in a magnetic field, using the full NLA equations and correct
    probabilities.

  \item \textbf{Analogue circuit design.}  Op-amp filter circuits
    designed by Claude, built on a breadboard, and tested---they worked.

  \item \textbf{Lecture notes from video.}  An entire 13-lecture course
    auto-transcribed and typeset as a 43-page PDF with every equation
    checked.
\end{enumerate}

\subsection{The FeO computation}

Returning to the live demo: by the end of the talk, Claude Code had
autonomously set up PySCF, written a production-quality computation
script, and begun calculating the potential energy curves of FeO\@.  The
workflow it produced:
\begin{itemize}
  \item SA-CASSCF(12,12) active space for 35 bond lengths.
  \item SC-NEVPT2 correlation corrections with CASCI intermediates (a
    subtlety it diagnosed itself when NEVPT2 on top of SA-CASSCF failed).
  \item X2C relativistic Hamiltonian with cc-pVTZ-DK basis set.
  \item Approximately 25 electronic states across multiple spin
    multiplicities.
\end{itemize}

The audience member who posed the challenge had left by the time results
appeared, but the computation was running and producing physically
reasonable output.  This illustrates both the power and the limitation:
the agent can \emph{execute}, but only a domain expert can
\emph{evaluate}.

\subsection{Highlights from the Q\&A}

The discussion after the talk covered several themes that are worth
recording.

\subsubsection{Literature reviews}

\begin{audience}[Q: Can agents do literature reviews?]
  Definitely, but they hallucinate references.  The guard rail: force the
  agent to \emph{download} every paper it cites to a local folder.
  String-match the filenames against the citations.  If the agent claims
  a paper exists that it has not downloaded, reject the citation.  This
  converts a hallucination-prone task into one with a verifiable paper
  trail.

  For training: take a paper you are interested in, do a reverse
  citation search, read through the results (bucket them---some are
  garbage, some are not), and iterate.  Once you have done this two or
  three times, it becomes an algorithm.  Then automate it.
\end{audience}

\subsubsection{Security and data privacy}

\begin{audience}[Q: Does the data I upload get remembered and used?]
  You have to assume it does.  Anthropic's business model depends on
  companies trusting that their IP is not leaked, and a single provable
  instance of data leaking would be an existential threat to their
  value as a company.  But ultimately this is a trust decision, the same
  kind you already make when you carry a phone in your pocket.

  Practical advice: run agents on a machine you do not care about.  A
  Mac Mini with nothing personal on it is the standard practice.  If it
  has arbitrary code execution capabilities, it can do arbitrary code
  execution.  There is no technical mitigation for that.
\end{audience}

\subsubsection{Reasoning capabilities}

\begin{audience}[Q: How is mathematical reasoning compatible with the
  model of a nondeterministic neural network?]
  Two factors.  First, scaling: a bigger neural network with in-context
  learning gets you $90\%$ of the way to reasoning-like capabilities.
  Second, behind the scenes, frontier labs run multi-agent critique:
  a proposer generates an argument, a critic finds flaws, and the
  result is iterated.  LLMs are only moderately good at reasoning, but
  they are much better at finding mistakes.  Combining the two gives an
  effective probabilistic oracle for correctness.
\end{audience}

\subsubsection{Correctness and verification}

\begin{audience}[Q: How do I know the agent understands what I am
  telling it?]
  There are several levels of correctness:
  \begin{enumerate}[(1)]
    \item \textbf{Formal proof}: every line is checked, the derivations
      are valid.  This can be done programmatically (e.g.\ Lean~4).
    \item \textbf{Compilation}: the code compiles, so it is at least
      syntactically correct.
    \item \textbf{Tests}: test-driven development is the primary
      mechanism.  Write the tests first (or have the agent write them),
      inspect the tests yourself as the source of truth, and then let
      the agent run them after every change.
    \item \textbf{Consensus}: peer review, replication, expert
      evaluation.
  \end{enumerate}
  The agent responds best when it has tests.  The magic is not in skills
  or special configuration---just tell it in natural language what you
  want and insist on tests.
\end{audience}

\subsubsection{Publication and formalisation}

\begin{audience}[Q: Should papers include formalisations?]
  We have entered a reproducibility crisis of epic proportions.  People
  can now one-shot a bad paper on the paid model of ChatGPT\@.  The
  peer-review system cannot cope with the load.

  If you can auto-formalise relatively painlessly, why not make it a
  requirement of publication?  The referee no longer needs to verify
  mathematical correctness---just semantic correctness and scientific
  interest.  For anything that uses linear algebra or functional
  analysis, there is no reason not to formalise it now.
\end{audience}

\subsection{Summary}

The entire talk built one idea from first principles:

\begin{keyresult}[The stack]
  \begin{center}
  \begin{tabular}{r@{\;$\to$\;}l}
    String in, string out & the LLM primitive \\
    Re-send history each turn & the illusion of chat \\
    Parse tool calls in a loop & the coding agent \\
    Multiple agents, one supervisor & multi-agent orchestration \\
    Human domain expertise & the source of ground truth \\
  \end{tabular}
  \end{center}
  Every component is simple and composable.  The gap between ``it
  hallucinated'' and ``it is a force multiplier'' is closed by
  technique and domain expertise, not by better models.
\end{keyresult}
